\section{Problem Formulation \& Environment}
A two-player zero-sum stochastic game is defined as a tuple $\langle S, A, P, R, \gamma \rangle$, where $S$ is the state space, $A = A_1 \times A_2$ is the joint action space, $P$ is the transition probability function, $R$ is the reward function, and $\gamma$ is the discount factor.

In this specific setting~\cite{littman1994markov}, the environment is a $4 \times 5$ gridworld soccer game where two agents ($A$ and $B$) compete to score goals. The state space $S$ consists of the distinct grid coordinates for both players and the binary ball possession status. Each player chooses from the individual action set $A_i = \{\text{North}, \text{South}, \text{East}, \text{West}, \text{Stand}\}$. The transition function $P$ governs simultaneous moves executed in uniformly random order, transferring ball possession without displacement whenever a player moves into the square occupied by the opponent. The reward function $R$ assigns $+1$ to the scoring player and $-1$ to the opponent, resetting the game afterward. The discount factor $\gamma = 0.9$ is emulated via a per-step draw termination probability of $1 - \gamma = 0.1$ (yielding $0$ reward to both players).

The core challenge stems from multi-agent non-stationarity: each agent must anticipate the opponent's evolving behavior while optimizing its own policy to maximize expected cumulative reward.
